\section{以西结- 「上帝坚固」}
http://www.wccec.org/newWCCEC/Version1/gym/marathon/Jeremiah.pdf
\subsection{以西結蒙召（1-3）}

\subsubsection{以西結蒙召做先知}
\begin{table}[H]
\centering
\begin{tabular}{p{1.2cm}|p{1.8cm}|p{13.cm}}\hline\hline
\multicolumn{1}{c|}{经文}&\multicolumn{1}{c|}{蒙召地方}&\multicolumn{1}{c}{内容}\\\hline

2:3&迦巴魯河邊&他對我說：「人子啊，\textcolor{red}{我差你往}悖逆的國民以色列人那裡去。他們是悖逆我的，他們和他們的列祖違背我直到今日。 \\\hline 
2:4&迦巴魯河邊&這眾子面無羞恥，心裡剛硬，\textcolor{red}{我差你往}他們那裡去，你要對他們說：『主耶和華如此說。』\\\hline 
2:7&迦巴魯河邊&他們或聽或不聽，你只管\textcolor{red}{將我的話告訴}他們，他們是極其悖逆的。\\\hline 
3:1&迦巴魯河邊&他對我說：「人子啊，要吃你所得的，要吃這書卷，好\textcolor{red}{去對以色列家講說}。」\\\hline 
3:4-5&迦巴魯河邊&他對我說：「人子啊，你往以色列家那裡去，\textcolor{red}{將我的話對他們講說}。 \textcolor{red}{你奉差遣}不是往那說話深奧、言語難懂的民那裡去，乃是往以色列家去。\\\hline
3:10-11&迦巴魯河邊&他又對我說：「人子啊，我對你所說的一切話，要心裡領會，耳中聽聞。你往你本國被擄的子民那裡去，他們或聽或不聽，你要對他們講說，告訴他們：『這是主耶和華說的。』」\\\hline
3:26-27&平原&我必使你的舌頭貼住上膛，以致你啞口，不能做責備他們的人，他們原是悖逆之家。 27 但我對你說話的時候，必使你開口，你就要對他們說：『主耶和華如此說。』聽的可以聽，不聽的任他不聽，因為他們是悖逆之家。\\\hline
\end{tabular}
\end{table}

\subsubsection{以西結蒙召做以色列家守望的人}
\begin{table}[H]
\centering
\begin{tabular}{p{1.2cm}|p{4.cm}|p{10.cm}}\hline\hline
\multicolumn{1}{c|}{经文}&\multicolumn{1}{c|}{蒙召时间}&\multicolumn{1}{c}{内容}\\\hline
3:16&五年四月初五日&人子啊，我立你做以色列家守望的人，所以你要聽我口中的話，替我警戒他們。\\\hline

32:17 /33:7&十二年十二月十五日&人子啊，我照樣立你做以色列家守望的人。所以你要聽我口中的話，替我警戒他們。 \\\hline
\end{tabular}
\end{table}

\subsubsection{以西結、以賽亞、耶利米蒙召对比}
\begin{table}[H]
\centering
\begin{tabular}{p{1.2cm}|p{2.cm}|p{2.cm}|p{10.cm}}\hline\hline
\multicolumn{1}{c|}{先知}&\multicolumn{1}{c|}{经文}&\multicolumn{1}{c|}{蒙召地方}&\multicolumn{1}{c}{内容}\\\hline
以賽亞&賽6:1-13&聖殿&看見耶和華在聖殿的寶座上，然後集中描述撒拉弗所代表和彰顯神的榮耀,從神在聖殿中榮耀的異象裏，領受他的使命\\\hline
耶利米&1:11-15&&年輕時所獲得的異象，展示了他使命的性質\\\hline 
以西結&結1:1-3:21&巴比倫&神向先知啟示祂的榮耀時，亦顯示了先知使命的性質\\\hline
\end{tabular}
\end{table}


\begin{table}[H]
\centering
\begin{tabular}{p{2.cm}|p{3.cm}|p{10.50cm}}\hline\hline
\multicolumn{1}{c|}{先知}&\multicolumn{1}{c|}{经文}&\multicolumn{1}{c}{内容}\\\hline
蒙召異象&1:28&以狂風開始，一朵大雲從北方趨近，雲的周圍有光，還有四活物和四個輪。活物頭上的寶座。在寶座上的，「就是耶和華榮耀的形象」\\\hline
服事对象&2:3&人子啊，我差你往悖逆的國民以色列人那裏去。\\\hline 
職事&3:11&斥責他那些被擄的同胞\\\hline
&3:16-17/33:1-9&成為以色列家的守望者\\\hline
&3:18-19&令整個國家悔改\\\hline
装备&2:9-10/3:1-3&吃一卷充滿哀號和悲痛的書卷，吃下時卻覺得甘甜如蜜\\\hline
蒙召的态度&2:6-7&對悖逆的以色列民產生恐懼憂悶消沈\\\hline
&3:9&他們雖是悖逆之家，你不要怕他們，也不要因他們的臉色驚惶\\\hline
&3:15&我就來到提勒亞畢——住在迦巴魯河邊被擄的人那裡，到他們所住的地方，在他們中間憂憂悶悶地坐了七日。\\\hline
事奉地点&3:25&被神限制在他的家中，只向那些尋求神的話和旨意的以色列人進行\\\hline
事奉原則&3:27/3:11&聽的可以聽，不聽的任他不聽，因為他們是悖逆之家\\\hline
事奉方法&20:49&会说比喻的人\\\hline
神捆綁&4:7-8&你要露出膀臂，面向被困的耶路撒冷，說預言攻擊這城。 8 我用繩索捆綁你，使你不能輾轉，直等你滿了困城的日子。\\\hline
人捆綁&3:24-25&耶和華對我說：「你進房屋去，將門關上。 25 人子啊，人必用繩索捆綁你，你就不能出去在他們中間來往。 \\\hline

啞口開口&3:26-27&我必使你的舌頭貼住上膛，以致你啞口，不能做責備他們的人，他們原是悖逆之家。 27 但我對你說話的時候，必使你開口，你就要對他們說：『主耶和華如此說。』\\\hline
\end{tabular}
\end{table}



\subsection{被擄中的信息}
\subsubsection{知道耶和華是神}
\begin{table}[H]
\centering
\begin{tabular}{p{1.5cm}|p{1.cm}|p{14.0cm}}\hline\hline
\multicolumn{1}{c|}{对象}&\multicolumn{1}{c|}{经文}&\multicolumn{1}{c}{内容}\\\hline

被掳犹大人&2:4-5&眾子面無羞恥，心裡剛硬，我差你往他們那裡去，你要對他們說：『主耶和華如此說。』 5 他們或聽或不聽（他們是悖逆之家），\textcolor{red}{必知道在他們中間有了先知}。\\\hline


亡国前的人&5:13&我要這樣成就怒中所定的。我向他們發的憤怒止息了，自己就得著安慰。我在他們身上成就怒中所定的，那時\textcolor{red}{他們就知道我耶和華}所說的是出於熱心。\\\hline
以色列人&6:5-7&我也要將以色列人的屍首放在他們的偶像面前，將你們的骸骨拋散在你們祭壇的四圍。...被杀的人必横尸你们中间，你们就知道我是耶和华。\\\hline
犹大人&6:10&\textcolor{red}{他们必知道我是耶和华}；我说要使这灾祸临到他们身上，并非空话。\\\hline

南国犹大&6:13&他們被殺的人倒在他們祭壇四圍的偶像中，就是各高岡、各山頂、各青翠樹下、各茂密的橡樹下，乃是他們獻馨香的祭牲給一切偶像的地方，那時\textcolor{red}{他們就知道我是耶和華}。 14 我必伸手攻擊他們，使他們的地從曠野到第伯拉他，一切住處極其荒涼，\textcolor{red}{他們就知道我是耶和華}。\\\hline

南国犹大&7:4&我必不顾惜你，也不怜悯你，我必按你的恶行报应你。这样，\textcolor{red}{你必知道我是耶和华}。\\\hline

南国犹大&7:9&我必不顾惜你，也不怜悯你，我必按你的恶行报应你。这样，\textcolor{red}{你就知道惩罚你的是我耶和华}。\\\hline
南国犹大&7:27&君要悲哀，王要披淒涼為衣，國民的手都發顫。我必照他們的行為待他們，按他們應得的審判他們，\textcolor{red}{他們就知道我是耶和華}。\\\hline
以色列&11:10&使你们死于刀下。我要审判整个以色列，这样\textcolor{red}{你们就知道我是耶和华}。\\\hline
以色列&11:12&这样\textcolor{red}{你们就知道我是耶和华}。因为你们没有遵行我的律例，没有顺从我的典章，却效法你们四邻外族的规矩。\\\hline
列國&12:15&我將他們四散在列國、分散在列邦的時候，\textcolor{red}{他們就知道我是耶和華}。 16 我卻要留下他們幾個人得免刀劍、饑荒、瘟疫，使他們在所到的各國中述說他們一切可憎的事，\textcolor{red}{人就知道我是耶和華}。\\\hline

以色列&12:20&有居民的城邑必變為荒場，地也必變為荒廢，\textcolor{red}{你們就知道我是耶和華}。\\\hline
以色列&13:9&我必动手攻击那些见假异象、说假预言的先知。他们必从我的子民当中被剔除，不被收录在以色列人的名册上，也不能进入以色列的土地。这样，\textcolor{red}{你们就知道我是主耶和华}。\\\hline
以色列&13:14&我必把你们粉饰的墙夷为平地，使其露出根基。墙倒塌的时候，你们必葬身其中。这样，\textcolor{red}{你们就知道我是耶和华}。\\\hline
發預言的女子&13:21&我也必撕裂你們下垂的頭巾，救我百姓脫離你們的手，不再被獵取落在你們手中，\textcolor{red}{你們就知道我是耶和華}。\\\hline

以色列家&13:23&因此，你们必不再见假异象，也不再占卜。我必从你们手中拯救我的子民。这样，\textcolor{red}{你们就知道我是耶和华}。\\\hline
以色列家&14:7-8&因為以色列家的人，或在以色列中寄居的外人，凡與我隔絕，將他的假神接到心裡，把陷於罪的絆腳石放在面前，又就了先知來要為自己的事求問我的，我耶和華必親自回答他。 8 我必向那人變臉，使他做了警戒、笑談，令人驚駭，並且我要將他從我民中剪除，\textcolor{red}{你們就知道我是耶和華}。\\\hline


被掳之人&14:22-23&然而，其中仍有幸免的人带着儿女逃出来。他们会逃到你们这里，你们看见他们的所作所为，便会对我在耶路撒冷所降的一切灾难感到欣慰。你们看见他们的所作所为，\textcolor{red}{就知道}我不是无缘无故地惩罚耶路撒冷。’这是主耶和华说的。\\\hline



以色列&15:7&我必向他們變臉，他們雖從火中出來，火卻要燒滅他們。我向他們變臉的時候，\textcolor{red}{你們就知道我是耶和華}\\\hline

西底家&17:21&他的一切軍隊，凡逃跑的都必倒在刀下，所剩下的也必分散四方，\textcolor{red}{你們就知道}說這話的是我耶和華。\\\hline

以色列人&20:12&又將我的安息日賜給他們好在我與他們中間為證據,\textcolor{red}{使他們知道我耶和華}是叫他們成為聖的.\\\hline
以色列人&20:20&且以我的安息日為聖。這日在我與你們中間為證據，\textcolor{red}{使你們知道我是耶和華你們的神}。
\\\hline

以色列人&20:26&因他們將一切頭生的經火，我就任憑他們在這供獻的事上玷汙自己，好叫他們淒涼，\textcolor{red}{使他們知道我是耶和華}。\\\hline


以色列人&20:38&我必從你們中間除淨叛逆和得罪我的人，將他們從所寄居的地方領出來，他們卻不得入以色列地，\textcolor{red}{你們就知道我是耶和華}。\\\hline

\end{tabular}
\end{table}


\begin{table}[H]
\centering
\begin{tabular}{p{1.5cm}|p{1.cm}|p{13.0cm}}\hline\hline
\multicolumn{1}{c|}{对象}&\multicolumn{1}{c|}{经文}&\multicolumn{1}{c}{内容}\\\hline

以色列人&20:42&我領你們進入以色列地，就是我起誓應許賜給你們列祖之地，那時\textcolor{red}{你們就知道我是耶和華}。\\\hline


以色列人&20:44&主耶和華說：以色列家啊，我為我名的緣故，不照著你們的惡行和你們的壞事待你們，\textcolor{red}{你們就知道我是耶和華}。\\\hline

有血氣的&20:48&\textcolor{red}{凡有血氣的都必知道}，是我耶和華使火著起，這火必不熄滅。\\\hline
以色列人&22:14-16&我耶和華說了這話，就必照著行。 15 我必將你分散在列國，四散在列邦，我也必從你中間除掉你的汙穢。 16 你必在列國人的眼前因自己所行的被褻瀆，\textcolor{red}{你就知道我是耶和華}。\\\hline

以色列人&22:20/22&人怎樣將銀、銅、鐵、鉛、錫聚在爐中，吹火熔化，照樣，我也要發怒氣和憤怒將你們聚集，放在城中熔化你們。銀子怎樣熔化在爐中，你們也必照樣熔化在城中，\textcolor{red}{你們就知道我耶和華}是將憤怒倒在你們身上了。\\\hline
以色列人& 23:49&人必照著你們的淫行報應你們，你們要擔當拜偶像的罪，\textcolor{red}{就知道我是主耶和華}。\\\hline

以色列人&24:24&以西結必這樣為你們做預兆，凡他所行的，你們也必照樣行。那事來到，\textcolor{red}{你們就知道我是主耶和華}。\\\hline

被掳以色列人&24:27&你必向逃脫的人開口說話，不再啞口。你必這樣為他們做預兆，\textcolor{red}{他們就知道我是耶和華}。\\\hline

亚扪人&25:5&我必使拉巴為駱駝場，使亞捫人的地為羊群躺臥之處，\textcolor{red}{你們就知道我是耶和華}。\\\hline
亚扪人&25:7&所以我伸手攻擊你，將你交給列國作為擄物。我必從萬民中剪除你，使你從萬國中敗亡。我必除滅你，\textcolor{red}{你就知道我是耶和華}。\\\hline

摩押人&25:11&我必向摩押施行審判，\textcolor{red}{他們就知道我是耶和華}。\\\hline
以東&25:14&我必藉我民以色列的手報復以東，以色列民必照我的怒氣，按我的憤怒在以東施報，\textcolor{red}{以東人就知道是我}施報。這是主耶和華說的。\\\hline
非利士人&25:17&我向他們大施報應，發怒斥責他們。我報復他們的時候，\textcolor{red}{他們就知道我是耶和華}。\\\hline
推羅居民&26:6&屬推羅城邑的居民必被刀劍殺滅，\textcolor{red}{他們就知道我是耶和華}。\\\hline

西頓&28:22-24&主耶和華如此說：西頓哪，我與你為敵，我必在你中間得榮耀。我在你中間施行審判，顯為聖的時候，\textcolor{red}{人就知道我是耶和華}。 23 我必使瘟疫進入西頓，使血流在她街上。被殺的必在其中仆倒，四圍有刀劍臨到她，\textcolor{red}{人就知道我是耶和華}。 24 四圍恨惡以色列家的人，必不再向他們做刺人的荊棘、傷人的蒺藜，\textcolor{red}{人就知道我是主耶和華}。\\\hline 


以色列家&28:25-26&主耶和華如此說：我將分散在萬民中的以色列家招聚回來，向他們在列邦人眼前顯為聖的時候，他們就在我賜給我僕人雅各之地仍然居住。 26 他們要在這地上安然居住。我向四圍恨惡他們的眾人施行審判以後，他們要蓋造房屋，栽種葡萄園，安然居住，就\textcolor{red}{知道我是耶和華他們的神}。\\\hline

埃及居民&29:6-7&埃及一切的居民，因向以色列家成了蘆葦的杖，\textcolor{red}{就知道我是耶和華}。 7 他們用手持住你，你就斷折，傷了他們的肩；他們倚靠你，你就斷折，閃了他們的腰。\\\hline

埃及居民&29:8-9&所以主耶和華如此說：我必使刀劍臨到你，從你中間將人與牲畜剪除。 9 埃及地必荒廢淒涼，\textcolor{red}{他們就知道我是耶和華}。因為法老說「這河是我的，是我所造的」\\\hline
以色列&29:16&埃及必不再做以色列家所倚靠的。以色列家仰望埃及人的時候，便思念罪孽，\textcolor{red}{他們就知道我是主耶和華}。\\\hline
以色列&29:21&當那日，我必使以色列家的角發生，又必使你以西結在他們中間得以開口，\textcolor{red}{他們就知道我是耶和華}。\\\hline

埃及人&30:8&我在埃及中使火著起，幫助埃及的都被滅絕，那時\textcolor{red}{他們就知道我是耶和華}\\\hline 

埃及人&30:19&我必这样审判埃及，\textcolor{red}{他們就知道我是耶和華}。\\\hline
埃及人&30:25-26&我必扶持巴比倫王的膀臂，法老的膀臂卻要下垂。我將我的刀交在巴比倫王手中，他必舉刀攻擊埃及地，\textcolor{red}{他們就知道我是耶和華}。我必將埃及人分散在列國，四散在列邦，\textcolor{red}{他們就知道我是耶和華}。\\\hline 
埃及人&32:15&我使埃及地變為荒廢淒涼，這地缺少從前所充滿的，又擊殺其中一切的居民，那時\textcolor{red}{他們就知道我是耶和華}。\\\hline

以色列&33:29&我因他們所行一切可憎的事，使地荒涼，令人驚駭，那時\textcolor{red}{他們就知道我是耶和華}。\\\hline
\end{tabular}
\end{table}


\begin{table}[H]
\centering
\begin{tabular}{p{1.5cm}|p{1.cm}|p{13.0cm}}\hline\hline
\multicolumn{1}{c|}{对象}&\multicolumn{1}{c|}{经文}&\multicolumn{1}{c}{内容}\\\hline
以色列&33:33&他們看你如善於奏樂、聲音幽雅之人所唱的雅歌，他們聽你的話卻不去行。 33 看哪，所說的快要應驗。應驗了，\textcolor{red}{他們就知道在他們中間有了先知}。\\\hline

以色列&34:27&田野的樹必結果，地也必有出產，他們必在故土安然居住。我折斷他們所負的軛，救他們脫離那以他們為奴之人的手，那時\textcolor{red}{他們就知道我是耶和華}。\\\hline

以色列&34:29-30&我必給他們興起有名的植物，他們在境內不再為饑荒所滅，也不再受外邦人的羞辱\textcolor{red}{必知道我耶和華他們的神是與他們同在}，\textcolor{red}{並知道他們以色列家是我的民}。這是主耶和華說的。\\\hline


西珥山&35:4&我必使你的城邑變為荒場，成為淒涼，\textcolor{red}{你就知道我是耶和華}。\\\hline

西珥山&35:9&我必使你永遠荒涼，使你的城邑無人居住，\textcolor{red}{你的民就知道我是耶和華}\\\hline

西珥山&35:12&\textcolor{red}{你也必知道我耶和華聽見了}你的一切毀謗，就是你攻擊以色列山的話說：『這些山荒涼，是歸我們吞滅的。』。\\\hline

以东&35:15&你既因以色列的荒凉而幸灾乐祸，我要让你遭受同样的下场。西珥山啊，你和以东都要荒凉。这样，\textcolor{red}{你们就知道我是耶和华}。\\\hline 


以色列家&36:11&我必使人和牲畜在你上面加增，他們必生養眾多。我要使你照舊有人居住，並要賜福於你比先前更多，\textcolor{red}{你就知道我是耶和華}。\\\hline 


以色列家&36:23&我要使我的大名顯為聖，這名在列國中已被褻瀆，就是你們在他們中間所褻瀆的。我在他們眼前，在你們身上顯為聖的時候，\textcolor{red}{他們就知道我是耶和華}。這是主耶和華說的。\\\hline 

以色列家&36:32&主耶和華說：\textcolor{red}{你們要知道}，我這樣行不是為你們。以色列家啊，當為自己的行為抱愧蒙羞！\\\hline 


外邦人&36:36&那時，\textcolor{red}{在你們四圍其餘的外邦人必知道我耶和華修造那毀壞之處}，培植那荒廢之地。我耶和華說過，也必成就。\\\hline 

&36:38&耶路撒冷在守節做祭物所獻的羊群怎樣多，照樣，荒涼的城邑必被人群充滿，\textcolor{red}{他們就知道我是耶和華}。\\\hline 



以色列全家&37:6&我必給你們加上筋，使你們長肉，又將皮遮蔽你們，使氣息進入你們裡面，你們就要活了，\textcolor{red}{你們便知道我是耶和華}。\\\hline

我的民&37:13-14&我的民哪，我開你們的墳墓，使你們從墳墓中出來，\textcolor{red}{你們就知道我是耶和華}。 14 我必將我的靈放在你們裡面，你們就要活了。我將你們安置在本地，\textcolor{red}{你們就知道我耶和華如此說，也如此成就了}。這是耶和華說的。\\\hline 


外邦人&37:28&我的聖所在以色列人中間直到永遠，\textcolor{red}{外邦人就必知道我是叫以色列成為聖的耶和華}。\\\hline 

各国的人&38:16&歌革啊，你必上來攻擊我的民以色列，如密雲遮蓋地面。末後的日子，我必帶你來攻擊我的地，到我在外邦人眼前，在你身上顯為聖的時候，\textcolor{red}{好叫他們認識我}。\\\hline 

各国的人&38:23&我必顯為大，顯為聖，在多國人的眼前顯現，\textcolor{red}{他們就知道我是耶和華}。\\\hline 

瑪各和海島&39:6&我要降火在瑪各和海島安然居住的人身上，他們就知道我是耶和華。\\\hline 

列國人&39:7&我要在我民以色列中顯出我的聖名，也不容我的聖名再被褻瀆，列國人就知道我是耶和華，以色列中的聖者。\\\hline 

各国的人/以色列家&39:21-23&我必顯我的榮耀在列國中，\textcolor{red}{萬民就必看見}我所行的審判與我在他們身上所加的手。 22 這樣，從那日以後，以色列家\textcolor{red}{必知道}我是耶和華他們的神。 23 列國人也\textcolor{red}{必知道}以色列家被擄掠，是因他們的罪孽。他們得罪我，我就掩面不顧，將他們交在敵人手中，他們便都倒在刀下。\\\hline 

以色列人&39:27-28&因我使他們被擄到外邦人中，後又聚集他們歸回本地，\textcolor{red}{他們就知道我是耶和華他們的神}。我必不再留他們一人在外邦。\\\hline 
\end{tabular}
\end{table}

\subsubsubsection{出埃及记}
\begin{table}[H]
\centering
\begin{tabular}{p{1.5cm}|p{1.cm}|p{13.0cm}}\hline\hline
\multicolumn{1}{c|}{对象}&\multicolumn{1}{c|}{经文}&\multicolumn{1}{c}{内容}\\\hline
埃及人&7:5&我伸手攻擊埃及，將以色列人從他們中間領出來的時候，埃及人就要知道我是耶和華。\\\hline
埃及人&14:4&我要使法老的心刚硬，使他派兵追赶你们。我要在法老和埃及军兵身上得到荣耀，好叫埃及人知道我是耶和华。”于是，以色列人依言而行。\\\hline
埃及人&14:18&当我在法老及其战车和骑兵身上得到荣耀时，埃及人就知道我是耶和华。\\\hline 
\end{tabular}
\end{table}


\subsubsubsection{列王记上下}
\begin{table}[H]
\centering
\begin{tabular}{p{1.5cm}|p{1.cm}|p{13.0cm}}\hline\hline
\multicolumn{1}{c|}{对象}&\multicolumn{1}{c|}{经文}&\multicolumn{1}{c}{内容}\\\hline
以色列人&上18:25-45&到了献晚祭的时候、先知以利亚近前来、说、亚伯拉罕、以撒、以色列的　神、耶和华阿、求你今日\textcolor{red}{使人知道你是以色列的神}、也\textcolor{red}{知道我是你的仆人}、又是奉你的命行这一切事。耶和华阿、求你应允我、应允我、\textcolor{red}{使这民知道你耶和华是神}、又\textcolor{red}{知道是你}叫这民的心回转。于是耶和华降下火来、烧尽燔祭、木柴、石头、尘土、又烧干沟里的水。众民看见了、就俯伏在地、说、耶和华是神、耶和华是神。\\\hline

亚哈&上 20:13&有一个先知来见以色列王亚哈说、耶和华如此说、这一大群人你看见了么．今日我必将他们交在你手里、你就知道我是耶和华。\\\hline 
天下万国&下19:19&我们的上帝耶和华啊，现在求你从亚述王手中拯救我们，让天下万国都知道唯有你是耶和华。\\\hline
\end{tabular}
\end{table}





\subsubsubsection{以賽亞書}
\begin{table}[H]
\centering
\begin{tabular}{p{1.5cm}|p{1.cm}|p{13.0cm}}\hline\hline
\multicolumn{1}{c|}{对象}&\multicolumn{1}{c|}{经文}&\multicolumn{1}{c}{内容}\\\hline
天下萬國&37:20&耶和華我們的神啊，現在求你救我們脫離亞述王的手，使天下萬國都知道，唯有你是耶和華。」\\\hline 
&66:19&我要顯神蹟在他們中間，逃脫的，我要差到列國去，就是到他施、普勒、拉弓的路德和土巴、雅完，並素來沒有聽見我名聲，沒有看見我榮耀遼遠的海島。他們必將我的榮耀傳揚在列國中。\\\hline 
\end{tabular}
\end{table}

\subsubsubsection{耶利米書 }
\begin{table}[H]
\centering
\begin{tabular}{p{1.5cm}|p{1.cm}|p{13.0cm}}\hline\hline
\multicolumn{1}{c|}{对象}&\multicolumn{1}{c|}{经文}&\multicolumn{1}{c}{内容}\\\hline
%%&14:18&我走到田间，眼前是具具被杀的尸体；我进入城里，见饥荒肆虐，祭司和先知也流亡异乡。\\\hline

&14:18&\\\hline
&16:21&耶和華說：「我要使他們知道，就是這一次使他們知道我的手和我的能力，他們\textcolor{red}{就知道我的名}是耶和華了。」\\\hline
犹大人&44:28&刀下余生、从埃及返回犹大的人必寥寥无几。那时，你们这些去埃及住的犹大人\textcolor{red}{必知道}谁的话会实现，是我的还是你们的。\\\hline 
\end{tabular}
\end{table}


\subsubsubsection{但以理書 }
\begin{table}[H]
\centering
\begin{tabular}{p{1.5cm}|p{1.cm}|p{13.0cm}}\hline\hline
\multicolumn{1}{c|}{对象}&\multicolumn{1}{c|}{经文}&\multicolumn{1}{c}{内容}\\\hline

尼布甲尼撒&4:17&這是守望者所發的命，聖者所出的令，好叫世人知道至高者在人的國中掌權，要將國賜予誰就賜予誰，或立極卑微的人執掌國權。\\\hline

尼布甲尼撒&4:25&你必被趕出離開世人，與野地的獸同居，吃草如牛，被天露滴濕，且要經過七期，等你知道至高者在人的國中掌權，要將國賜予誰就賜予誰。\\\hline

尼布甲尼撒&4:35-37&世人都微不足道，祂在天军和世人中独行其道，无人能拦阻祂的手，或质问祂的作为。“那时，我恢复了神智，也恢复了我的威严和荣耀，重现我国度的辉煌。我的谋士和大臣都来朝见我，我重掌国权，势力比以前更大。现在我尼布甲尼撒颂赞、尊崇、敬奉天上的王，因为祂的作为公正，祂行事公平，能够贬抑行为骄傲的人。”\\\hline

尼布甲尼撒&5:21&他被趕出離開世人，他的心變如獸心，與野驢同居，吃草如牛，身被天露滴濕，等他知道至高的神在人的國中掌權，憑自己的意旨立人治國。\\\hline
大利乌&6:26-27&我下令，我统治的国民都要敬畏但以理的上帝，“因为祂是永活长存的上帝，祂的国度永不灭亡，祂的统治直到永远。祂庇护、拯救，在天上地下行神迹奇事，救但以理脱离狮子的口。\\\hline 
\end{tabular}
\end{table}


\subsubsubsection{詩篇}
\begin{table}[H]
\centering
\begin{tabular}{p{5.5cm}|p{1.cm}|p{9.50cm}}\hline\hline
\multicolumn{1}{c|}{对象}&\multicolumn{1}{c|}{经文}&\multicolumn{1}{c}{内容}\\\hline
&9:16&耶和華已將自己顯明了，他已施行審判，惡人被自己手所做的纏住了。\\\hline
&59:136&求你發怒，使他們消滅，以至歸於無有，叫他們知道神在雅各中間掌權，直到地極。\\\hline 
以東,以實瑪利,摩押,夏甲人,迦巴勒,亞捫,亞瑪力,非利士,推羅,亞述&83:17-18&願他們永遠羞愧驚惶，願他們慚愧滅亡，使他們知道，唯獨你，名為耶和華的，是全地以上的至高者。\\\hline 

\end{tabular}
\end{table}



 
\subsubsection{審判的信息以色列的罪}
\begin{table}[H]
\centering
\begin{tabular}{p{3.75cm}|p{2.0cm}|p{11.0cm}}\hline\hline
\multicolumn{1}{c|}{審判}&\multicolumn{1}{c|}{经文}&\multicolumn{1}{c}{寓意}\\\hline
罪&5:6-7&神向以西結解釋以色列的背道\\\hline 
審判&5:8-13&審判將會持續，直至以色列人承認他們的主信實守約，向他們施行公義的審判\\\hline 
向以色列眾山說預言&6:1-10&以色列因拜偶像受罰\\\hline 
拍手頓足&6:11-14&一切住處極其荒涼\\\hline
一災、杖開花&7:1-22&全地都作可憎的惡事，導致神的審判要降在地上和百姓身上,期望他們再次承認祂是他們的神\\\hline 
要製造鎖鏈&7:23-27&以兆被擄\\\hline 
百姓的復興&11:17&先知在三十三至四十八章再論述這4個主題。\\\hline 
土地的復興&11:17&先知在三十三至四十八章再論述這4個主題。\\\hline 
百姓的潔淨&11:18&先知在三十三至四十八章再論述這4個主題。\\\hline 
神與祂的百姓重新建立相交&11:19-20&先知在三十三至四十八章再論述這4個主題。\\\hline 
澄清谣言&12:21-25&日子遲延，一切異象都落了空\\\hline 
&12:26-28&他所見的異象是關乎後來許多的日子，所說的預言是指著極遠的時候。\\\hline 
&18&父親吃了酸葡萄，兒子的牙酸倒了\\\hline 
家中信息&14:1-11&幾個以色列長老將他們的假神接到心裡\\\hline 
&14:12-23&四樣大災，就是刀劍、饑荒、惡獸、瘟疫降在耶路撒冷\\\hline 
&20:1-44&向幾個長老历数其列祖在埃及曠野迦南之違逆\\\hline 
\end{tabular}
\end{table}

\subsubsection{以色列对審判的信息错误看法}
\begin{table}[H]
\centering
\begin{tabular}{p{3.75cm}|p{2.0cm}|p{11.0cm}}\hline\hline
\multicolumn{1}{c|}{審判}&\multicolumn{1}{c|}{经文}&\multicolumn{1}{c}{寓意}\\\hline
異象落空&12:22&日子遲延，一切異象都落了空\\\hline 
極遠的時候&12:27&以色列家的人說：『他所見的異象是關乎後來許多的日子，所說的預言是指著極遠的時候。』\\\hline 
错误根源&13&這些假先知在百姓中煽動了邪惡、說謊和欺騙的行為\\\hline 
父债子还&18&你們在以色列地怎麼用這俗語說『父親吃了酸葡萄，兒子的牙酸倒了』\\\hline 
\end{tabular}
\end{table}
　　
\subsubsection{耶和華的\textcolor{red}{話}[又]臨到我}

\subsubsection{主的\textcolor{red}{靈}[或作手]在我身上}
\begin{table}[H]
\centering
\begin{tabular}{p{2.cm}|p{12.50cm}}\hline\hline
\multicolumn{1}{c|}{经文}&\multicolumn{1}{c}{内容}\\\hline
39-42章&「耶和華的話[又]臨到我\\\hline
\end{tabular}
\end{table}

\subsubsection{預言}
\begin{table}[H]
\centering
\begin{tabular}{p{2.cm}|p{12.50cm}}\hline\hline
\multicolumn{1}{c|}{经文}&\multicolumn{1}{c}{内容}\\\hline
37:27-28&大卫的后裔\\\hline
&以色列的復興\\\hline
&犹大的亡国\\\hline
38-39:20&瑪各地的歌革\\\hline
39:25-28&以色列蒙恩復歸故土\\\hline
39:29&將我的靈澆灌以色列家\\\hline
\end{tabular}
\end{table}



\subsection{比喻}
\begin{table}[H]
\centering
\begin{tabular}{p{1.75cm}|p{1.50cm}|p{13.0cm}}\hline\hline
\multicolumn{1}{c|}{比喻}&\multicolumn{1}{c|}{经文}&\multicolumn{1}{c}{寓意}\\\hline

荊棘蒺蠍子&2:6&人子啊，雖有荊棘和蒺藜在你那裡，你又住在蠍子中間，總不要怕他們，也不要怕他們的話。他們雖是悖逆之家，還不要怕他們的話，也不要因他們的臉色驚惶。\\\hline 
金鋼鑽火石&3:7-8&以色列家卻不肯聽從你，因為他們不肯聽從我，原來以色列全家是額堅心硬的人。 8 看哪，我使你的臉硬過他們的臉，使你的額硬過他們的額。 9 我使你的額像金鋼鑽，比火石更硬。\\\hline

荒場狐狸&13:1-7&以色列中說預言的先知沒有上去堵擋破口\\\hline
未泡透灰&13:8-16&假先知的預言如立起牆壁的必倒塌\\\hline
靠枕頭巾&13:17-23&向本民中從己心發預言的女子為要獵取人的性命\\\hline
焚葡萄枝& 15&以葡萄木被焚喻遭受徹底的毀滅\\\hline 
三個淫婦&16&妹妹\textcolor{red}{所多瑪}（創19）和姐姐\textcolor{red}{撒瑪利亞}（王下17:6）被稱為耶路撒冷的姊妹（結16:46）。這兩座城不及耶路撒冷的淫蕩（16:48-51）。耶路撒冷在悔改之後，將重新回復立約的福祉（16:62、63）\\\hline
两个淫妇，阿荷拉/阿荷利巴&23&兩個女子是一母所生. 她們在埃及行邪淫,在幼年時行邪淫。她們在那裡做處女的時候,有人擁抱她們的懷,撫摸她們的乳. 她們的名字，姐姐名叫阿荷拉就是撒馬利亞,妹妹名叫阿荷利巴是耶路撒冷。\\\hline
二鷹二葡萄樹&17&尼布甲尼撒是「一大鷹，翅膀大，翎毛長，羽毛豐滿，彩色俱備」，他將約雅斤「香柏樹梢」從王位中除去，將他的「香柏樹儘尖的嫩枝」──即猶大國的青年領袖一同擄去（17:3-4）。猶大（一棵葡萄樹）試圖與「又有一大鷹，翅膀大，羽毛多」（17:7）的埃及法老合弗拉結盟對抗尼布甲尼撒，導致尼布甲尼撒將葡萄樹連根拔起使它枯乾（17:9-10）。神應許在被擄之後，讓祂的百姓在彌賽亞「一嫩枝」（17:22）的帶領下，重回他們的國土。嫩枝象徵彌賽亞會長成一棵佳美的香柏樹，讓鳥兒得著蔭庇和保護。（「耶和華使高樹矮小，矮樹高大」，17:24）\\\hline

一只厉鹰&17:3&迦勒底 人是一只厉鹰\\\hline
香柏树&17:3-4&折去香柏樹儘尖的嫩枝，叼到貿易之地，放在買賣城中\\\hline
香柏高枝&17:22-24&將香柏樹梢擰去栽於以色列高處的山田野的樹木都必知道耶和華\\\hline
&31:3-5&埃及也是一株香柏树\\\hline
一只鳄鱼&32:1-3&埃及\\\hline
母獅小獅&19:1-9&那隻母獅是約西亞王的妻子哈慕他（王下23:31），生了兩個兒子，約哈斯和西底家。約哈斯在以西結書19:3-4中，描述為一隻長大後被帶到埃及的小獅（主前508年，被法老尼哥帶往埃及；王下23:31-34）。西底家在10年後繼承王位。在哀歌中，先知運用想象力，將小獅代表西底家，他最終成為背叛的統治者，而被帶到巴比倫（結19:7-9）\\\hline
葡萄樹&19:10-14&枝幹折斷枯乾，被火燒毀了。 13 如今，栽於曠野乾旱無水之地。 14 火也從它枝幹中發出，燒滅果子，以致沒有堅固的枝幹可做掌權者的杖。\\\hline
樹人見焚&20:45-49&南方的一切青樹和枯樹見焚，猛烈的火焰必不熄滅。從南到北，人的臉面都被燒焦。\\\hline
拔刀出鞘&21:1-32&自南至北攻擊一切有血氣的，以色列王及亞捫\\\hline
銅鐵鉛錫&22:17-31&\textcolor{red}{先知}同謀背叛，如咆哮的獅子抓撕掠物，他們吞滅人民，搶奪財寶，使這地多有寡婦。其中的\textcolor{red}{祭司}強解我的律法，褻瀆我的聖物，不分別聖的和俗的，也不使人分辨潔淨的和不潔淨的，又遮眼不顧我的安息日，我也在他們中間被褻慢。其中的\textcolor{red}{首領}彷彿豺狼抓撕掠物，殺人流血，傷害人命，要得不義之財。其中的\textcolor{red}{先知}為百姓用未泡透的灰抹牆，就是為他們見虛假的異象，用謊詐的占卜，說『主耶和華如此說』，其實耶和華沒有說。\textcolor{red}{ 國內眾民}一味地欺壓，慣行搶奪，虧負困苦窮乏的，背理欺壓寄居的。 \\\hline


鍋煮肥肉&24:1-14&『主耶和華如此說：將鍋放在火上，放好了，就倒水在其中。 4 將肉塊，就是一切肥美的肉塊，腿和肩都聚在其中，拿美好的骨頭把鍋裝滿。 5 取羊群中最好的，將柴堆在鍋下，使鍋開滾，好把骨頭煮在其中。\\\hline
牧人找羊&34&預言其僕必牧群羊\\\hline

\end{tabular}
\end{table}


\subsection{预兆}
\begin{table}[H]
\centering
\begin{tabular}{p{1.75cm}|p{1.50cm}|p{13.0cm}}\hline\hline
\multicolumn{1}{c|}{预兆}&\multicolumn{1}{c|}{经文}&\multicolumn{1}{c}{寓意}\\\hline
吃書卷&2:8-10&你要開口吃我所賜給你的。」 9 我觀看，見有一隻手向我伸出來，手中有一書卷。 10 他將書卷在我面前展開，內外都寫著字，其上所寫的有哀號、嘆息、悲痛的話。\\\hline
&3:1-3&他對我說：「人子啊，要吃你所得的，要吃這書卷，好去對以色列家講說。」 2 於是我開口，他就使我吃這書卷。 3 又對我說：「人子啊，要吃我所賜給你的這書卷，充滿你的肚腹。」我就吃了，口中覺得其甜如蜜。\\\hline

被人捆綁&3:24-25&耶和華對我說：「你進房屋去，將門關上。 25 人子啊，人必用繩索捆綁你，你就不能出去在他們中間來往。\\\hline

啞口&3:26-27&我必使你的舌頭貼住上膛，以致你啞口，不能做責備他們的人，他們原是悖逆之家。 27 但我對你說話的時候，必使你開口，你就要對他們說：『主耶和華如此說。』聽的可以聽，不聽的任他不聽，因為他們是悖逆之家。\\\hline
開口&24:25-27&那日逃脫的人豈不來到你這裡，使你耳聞這事嗎？ 27 你必向逃脫的人開口說話，不再啞口。你必這樣為他們做預兆，他們就知道我是耶和華。\\\hline
&29:21&當那日，我必使以色列家的角發生，又必使你以西結在他們中間得以開口，他們就知道我是耶和華。\\\hline
&33:21-22&我們被擄之後十二年十月初五日，有人從耶路撒冷逃到我這裡，說城已攻破。 22 逃來的人未到前一日的晚上，耶和華的靈降在我身上，開我的口。到第二日早晨，那人來到我這裡，我口就開了，不再緘默。\\\hline

一塊磚&4:1&人子啊，你要拿一塊磚，擺在你面前，將一座耶路撒冷城畫在其上\\\hline
撞錘攻城&4:2-3&又圍困這城，造臺築壘，安營攻擊，在四圍安設撞錘攻城。 3 又要拿個鐵鏊，放在你和城的中間作為鐵牆。你要對面攻擊這城，使城被困。\\\hline
&21:22&巴比倫王右手中拿著為耶路撒冷占卜的籤，使他安設撞城錘，張口叫殺，揚聲呐喊,築壘造臺,以撞城錘,攻打城門。\\\hline
左右側臥&4:4-8&以西結以側臥來代表以色列猶大家的罪孽\\\hline
按量吃喝&4:9-17&代表了耶路撒冷在被圍困的最後日子中的悲哀和恐懼\\\hline
頭髮被剃&5:1-17&代表了耶路撒冷的命運\\\hline
墙挖洞&12:1-16&把墙挖洞，肩带包袱，蒙住脸走出去\\\hline
吃飯喝水&12:17-20&你吃飯必膽戰，喝水必惶惶憂慮。\\\hline
喪妻勿哭&24:15-24&以西結勿為妻死哀哭喻以色列人受苦不暇悲傷\\\hline

二杖為一&37:15-28&預表以色列猶大必合為一\\\hline
\end{tabular}
\end{table}


\subsection{異象1/8:4/10:1-22/11:22-23/43:3}
\begin{table}[H]
\centering
\begin{tabular}{p{1.85cm}|p{2.0cm}|p{13.0cm}}\hline\hline
\multicolumn{1}{c|}{異象}&\multicolumn{1}{c|}{经文}&\multicolumn{1}{c}{寓意}\\\hline
狂风&1:4&从北方刮来一朵火云：\\\hline
四个活物&1:5-14/10：6-8,12-15&各有四个脸面、四个翅膀，又有腿和手。四活物移动时，不论朝那个方向都能够 往前直行，不必转身；他们的行动迅速，像电光一闪。\\\hline
四轮&1:15-21/10:9-17&四活物的脸旁各有一轮在地上，形状和颜色好像水苍玉，其 作法好像轮中套轮。它们移动时像四活物一样，不论朝那个方向都能够 往前直行，不必转身\\\hline
神的荣耀&1:22-28&\textcolor{red}{迦巴魯河邊}看到宝座彷佛蓝宝石；荣耀的形像彷佛人的形像，周围有光辉像彩虹\\\hline

&3:12-14&那時，靈將我舉起，我就聽見在我身後有震動轟轟的聲音說：「從耶和華的所在顯出來的榮耀是該稱頌的！」 13 我又聽見那活物翅膀相碰與活物旁邊輪子旋轉震動轟轟的響聲。於是靈將我舉起，帶我而去。我心中甚苦，靈性憤激，並且耶和華的靈在我身上大有能力。\\\hline

&3:22-24&耶和華的靈在那裡降在我身上，他對我說：「你起來，往平原去，我要在那裡和你說話。」 23 於是我起來往平原去，不料耶和華的榮耀正如我在迦巴魯河邊所見的一樣，停在那裡，我就俯伏於地。 24 靈就進入我裡面，使我站起來。\\\hline

神的榮耀&8:2-4&耶路撒冷\textcolor{red}{朝北的內院}有以色列神的榮耀，形狀與我在平原所見的一樣\\\hline

污辱圣殿&8:1-6&耶路撒冷朝北的內院門口有觸動主怒偶像的座位\\\hline
&8:7-13&四面牆上畫著各樣爬物和可憎的走獸，七十個長老拿香爐煙雲的香氣上騰\\\hline
&8:14-15&外院朝北的門口有婦女坐著為搭模斯哭泣\\\hline
&8:16-18&在耶和華殿門口，廊子和祭壇中間有二十五個人背向耶和華的殿，面向東方拜日頭\\\hline
神的榮耀&9:3&以色列神的榮耀本在基路伯上，現今從那裡\textcolor{red}{升到殿的門檻}。\\\hline
滅命者&9&從聖所起全都殺盡，只是凡有記號的人，不要挨近他。\\\hline

神的榮耀&10:1,3-5&基路伯站在\textcolor{red}{殿的右邊}，雲彩充滿了內院。 4 耶和華的榮耀從基路伯那裡上升，停在\textcolor{red}{門檻以上}。殿內滿了雲彩，院宇也被耶和華榮耀的光輝充滿。 5 基路伯翅膀的響聲聽到\textcolor{red}{外院}，好像全能神說話的聲音。\\\hline


&10:18-19&耶和華的榮耀從\textcolor{red}{殿的門檻那裡出去}，停在基路伯以上。 19 基路伯出去的時候，就展開翅膀，在我眼前離地上升，輪也在他們的旁邊，都\textcolor{red}{停在耶和華殿的東門口}。在他們以上有以色列神的榮耀。\\\hline
穿細麻衣人&10:1-6&從基路伯中間將火炭取滿兩手，撒在城上。\\\hline

民牧謀惡者&11:1-2&殿向東的東門，誰知，在門口有二十五個人,其中有民間的首領押朔的兒子雅撒尼亞和比拿雅的兒子毗拉提是圖謀罪孽的人，在這城中給人設惡謀\\\hline

奉命說審判&11:3-13&毗拉提死了，以西結為遺民禱主\\\hline

賜肉心新靈&11:14-21&從萬民中被招聚得以色列地，我要使他們有合一的心，也要將新靈放在他們裡面，又從他們肉體中除掉石心，賜給他們肉心\\\hline


神的榮耀&11:22-23&耶和華的榮耀從城中上升，停在\textcolor{red}{城東的那座山上}。\\\hline



枯骨復生&37:1-14&示以色列家絕望仍得復興\\\hline

圣殿圣城&40-48&感靈見以色列地有建城之異象\\\hline
\end{tabular}
\end{table}


\subsection{预言}
\begin{wrapfigure}{R}{0.6\textwidth}
\centering
\includegraphics[width=0.95\textwidth]{ZongJie/figs/YiXiJie/altar}
\end{wrapfigure}

\begin{wrapfigure}{R}{0.6\textwidth}
\centering
\includegraphics[width=0.95\textwidth]{ZongJie/figs/YiXiJie/fourBadVisions}
\end{wrapfigure}


\begin{wrapfigure}{R}{0.6\textwidth}
\centering
\includegraphics[width=0.95\textwidth]{ZongJie/figs/YiXiJie/FourVisions}
\end{wrapfigure}


\begin{wrapfigure}{R}{0.6\textwidth}
\centering
\includegraphics[width=0.95\textwidth]{ZongJie/figs/YiXiJie/Jerusalem1}
\end{wrapfigure}

\begin{wrapfigure}{R}{0.6\textwidth}
\centering
\includegraphics[width=0.95\textwidth]{ZongJie/figs/YiXiJie/Jerusalem2}
\end{wrapfigure}


\begin{wrapfigure}{R}{0.6\textwidth}
\centering
\includegraphics[width=0.95\textwidth]{ZongJie/figs/YiXiJie/kings}
\end{wrapfigure}


\begin{wrapfigure}{R}{0.6\textwidth}
\centering
\includegraphics[width=0.95\textwidth]{ZongJie/figs/YiXiJie/life}
\end{wrapfigure}


\begin{wrapfigure}{R}{0.6\textwidth}
\centering
\includegraphics[width=0.95\textwidth]{ZongJie/figs/YiXiJie/mainMessage}
\end{wrapfigure}


\begin{wrapfigure}{R}{0.6\textwidth}
\centering
\includegraphics[width=0.95\textwidth]{ZongJie/figs/YiXiJie/OuXiang1}
\end{wrapfigure}

\begin{wrapfigure}{R}{0.6\textwidth}
\centering
\includegraphics[width=0.95\textwidth]{ZongJie/figs/YiXiJie/paragraphs}
\end{wrapfigure}


\begin{wrapfigure}{R}{0.6\textwidth}
\centering
\includegraphics[width=0.95\textwidth]{ZongJie/figs/YiXiJie/Prest}
\end{wrapfigure}


\begin{wrapfigure}{R}{0.6\textwidth}
\centering
\includegraphics[width=0.95\textwidth]{ZongJie/figs/YiXiJie/sonOfMan}
\end{wrapfigure}

\begin{wrapfigure}{R}{0.6\textwidth}
\centering
\includegraphics[width=0.95\textwidth]{ZongJie/figs/YiXiJie/summery}
\end{wrapfigure}


\begin{wrapfigure}{R}{0.6\textwidth}
\centering
\includegraphics[width=0.95\textwidth]{ZongJie/figs/YiXiJie/Vision}
\end{wrapfigure}

\begin{wrapfigure}{R}{0.6\textwidth}
\centering
\includegraphics[width=0.95\textwidth]{ZongJie/figs/YiXiJie/Altars}
\end{wrapfigure}



%\begin{changemargin}{-1cm}{-1cm}
\begin{sidewaystable} % <-- HERE
\begin{tikzpicture}[thick,snake=zigzag, line before snake = 5mm, line after snake = 5mm]

     % draw 1st horizontal line 
     \draw (-1.,1.55) node[below=3pt] {} node[above=3pt] {以赛亚};
     %\draw (0.,2) node[below=3pt] {} node[above=3pt] {1};                         
     \draw (0.5,2) node[below=3pt] {} node[above=3pt] {};                         
    % \draw (1.,2) node[below=3pt] {} node[above=3pt] {3};                         
    % \draw (1.5,2) node[below=3pt] {} node[above=3pt] {4};                                        
     \draw (2.,2) node[below=3pt] {} node[above=3pt] {};                              
  %  \draw (2.5,2) node[below=3pt] {} node[above=3pt] {6};                  
  %  \draw (21.5,2) node[below=3pt] {} node[above=3pt] {7};                                                       
 
     \draw (0.5,4.4) node[below=3pt] {} node[above=3pt] {亚};   
    \draw (0.5,4.) node[below=3pt] {} node[above=3pt] {兰};  
     \draw (0.5,3.6) node[below=3pt] {} node[above=3pt] {被};   
    \draw (0.5,3.2) node[below=3pt] {} node[above=3pt] {亚};   
     \draw (0.5,2.8) node[below=3pt] {} node[above=3pt] {述};   
    \draw (0.5,2.4) node[below=3pt] {} node[above=3pt] {灭}; 
   
     \draw (1.5,4.4) node[below=3pt] {} node[above=3pt] {北};   
    \draw (1.5,4.) node[below=3pt] {} node[above=3pt] {国};  
     \draw (1.5,3.6) node[below=3pt] {} node[above=3pt] {被};   
    \draw (1.5,3.2) node[below=3pt] {} node[above=3pt] {亚};   
     \draw (1.5,2.8) node[below=3pt] {} node[above=3pt] {述};   
    \draw (1.5,2.4) node[below=3pt] {} node[above=3pt] {掳}; 
                               
     \draw (2.5,4.8) node[below=3pt] {} node[above=3pt] {南};   
    \draw (2.5,4.4) node[below=3pt] {} node[above=3pt] {国};                                
     \draw (2.5,4.0) node[below=3pt] {} node[above=3pt] {被};   
    \draw (2.5,3.6) node[below=3pt] {} node[above=3pt] {巴};                                
     \draw (2.5,3.2) node[below=3pt] {} node[above=3pt] {比};   
    \draw (2.5,2.8) node[below=3pt] {} node[above=3pt] {伦};                                    
    \draw (2.5,2.4) node[below=3pt] {} node[above=3pt] {掳};                                                   

  % \draw (3.5,3.6) node[below=3pt] {} node[above=3pt] {天使印};                                         
   \draw (3.5,3.2) node[below=3pt] {} node[above=3pt] {南};                                         
   \draw (3.5,2.8) node[below=3pt] {} node[above=3pt] {国};                                            
   \draw (3.5,2.4) node[below=3pt] {} node[above=3pt] {归};                                       
   \draw (3.5,2.0) node[below=3pt] {} node[above=3pt] {回};                                            


   \draw (6,3.6) node[below=3pt] {} node[above=3pt] {弥};                                         
   \draw (6,3.2) node[below=3pt] {} node[above=3pt] {赛};                                         
   \draw (6,2.8) node[below=3pt] {} node[above=3pt] {亚};                                            
   \draw (6,2.4) node[below=3pt] {} node[above=3pt] {一次};                                          
  \draw (6,2.0) node[below=3pt] {} node[above=3pt] {来};                                                
%  \draw (6,1.6) node[below=3pt] {} node[above=3pt] {弥};                                         

   \draw (7,3.2) node[below=3pt] {} node[above=3pt] {耶};                                         
   \draw (7,2.8) node[below=3pt] {} node[above=3pt] {稣};                                         
   \draw (7,2.4) node[below=3pt] {} node[above=3pt] {的};                                            
   \draw (7,2.) node[below=3pt] {} node[above=3pt] {工作};                                          
 % \draw (7,1.6) node[below=3pt] {} node[above=3pt] {死};                                                

   \draw (8,3.2) node[below=3pt] {} node[above=3pt] {耶};                                         
   \draw (8,2.8) node[below=3pt] {} node[above=3pt] {稣};                                         
   \draw (8,2.4) node[below=3pt] {} node[above=3pt] {被};                                            
   \draw (8,2.) node[below=3pt] {} node[above=3pt] {钉死};                                          
 
  \draw (9,3.2) node[below=3pt] {} node[above=3pt] {耶};                                         
   \draw (9,2.8) node[below=3pt] {} node[above=3pt] {稣};                                         
   \draw (9,2.4) node[below=3pt] {} node[above=3pt] {复};                                            
   \draw (9,2.) node[below=3pt] {} node[above=3pt] {活};                                          



   \draw (10,3.2) node[below=3pt] {} node[above=3pt] {圣};                                         
   \draw (10,2.8) node[below=3pt] {} node[above=3pt] {灵};                                         
   \draw (10,2.4) node[below=3pt] {} node[above=3pt] {降临};                                            
   \draw (10,2.) node[below=3pt] {} node[above=3pt] {临};                                          


  
   \draw (12,3.2) node[below=3pt] {} node[above=3pt] {教};                                         
   \draw (12,2.8) node[below=3pt] {} node[above=3pt] {会};                                            
   \draw (12,2.4) node[below=3pt] {} node[above=3pt] {被};                                          
   \draw (12,2.) node[below=3pt] {} node[above=3pt] {提};                                                



   \draw (15,3.6) node[below=3pt] {} node[above=3pt] {弥};                                         
   \draw (15,3.2) node[below=3pt] {} node[above=3pt] {赛};                                         
   \draw (15,2.8) node[below=3pt] {} node[above=3pt] {亚};                                            
   \draw (15,2.4) node[below=3pt] {} node[above=3pt] {二次};                                          
  \draw (15,2.) node[below=3pt] {} node[above=3pt] {来};                                                


   \draw (16,3.6) node[below=3pt] {} node[above=3pt] {以色};                                         
   \draw (16,3.2) node[below=3pt] {} node[above=3pt] {列和};                                         
   \draw (16,2.8) node[below=3pt] {} node[above=3pt] {犹大};                                            
   \draw (16,2.4) node[below=3pt] {} node[above=3pt] {二次};                                          
  \draw (16,2.) node[below=3pt] {} node[above=3pt] {归回};                                                


  \draw (17,3.2) node[below=3pt] {} node[above=3pt] {一};                                         
   \draw (17,2.8) node[below=3pt] {} node[above=3pt] {次};                                            
   \draw (17,2.4) node[below=3pt] {} node[above=3pt] {复};                                          
  \draw (17,2.) node[below=3pt] {} node[above=3pt] {活};                                                

  \draw (18,3.2) node[below=3pt] {} node[above=3pt] {审};                                         
   \draw (18,2.8) node[below=3pt] {} node[above=3pt] {判};                                            
   \draw (18,2.4) node[below=3pt] {} node[above=3pt] {列};                                          
  \draw (18,2.) node[below=3pt] {} node[above=3pt] {国};                                                


  \draw (19.2,3.2) node[below=3pt] {} node[above=3pt] {千};                                                
  \draw (19.2,2.8) node[below=3pt] {} node[above=3pt] {禧};                                                   \draw (19.2,2.4) node[below=3pt] {} node[above=3pt] {年};                                                    \draw (19.2,2.) node[below=3pt] {} node[above=3pt] {开始};                                                

  \draw (21.5,4.8) node[below=3pt] {} node[above=3pt] {撒旦};                                                
  \draw (21.5,4.4) node[below=3pt] {} node[above=3pt] {出};                                                   \draw (21.5,4.) node[below=3pt] {} node[above=3pt] {无底坑};                                                    \draw (21.5,3.6) node[below=3pt] {} node[above=3pt] {诱惑};                                                
\draw (21.5,3.2) node[below=3pt] {} node[above=3pt] {列国};                                               
\draw (21.5,2.8) node[below=3pt] {} node[above=3pt] {围攻};                                                    \draw (21.5,2.4) node[below=3pt] {} node[above=3pt] {耶路};                                                
\draw (21.5,2.) node[below=3pt] {} node[above=3pt] {撒冷};                                               


\draw (22.5,4.8) node[below=3pt] {} node[above=3pt] {所有};                                                
\draw (22.5,4.4) node[below=3pt] {} node[above=3pt] {的人};                                                   \draw (22.5,4.) node[below=3pt] {} node[above=3pt] {复活};                                                   
\draw (22.5,3.6) node[below=3pt] {} node[above=3pt] {白色};                                                
\draw (22.5,3.2) node[below=3pt] {} node[above=3pt] {大};                                                   \draw (22.5,2.8) node[below=3pt] {} node[above=3pt] {宝};                                                    \draw (22.5,2.4) node[below=3pt] {} node[above=3pt] {座};                                                
\draw (22.5,2.0) node[below=3pt] {} node[above=3pt] {审判};                                                

  \draw (23.5,3.2) node[below=3pt] {} node[above=3pt] {进};                                                
  \draw (23.5,2.8) node[below=3pt] {} node[above=3pt] {入};                                                
    \draw (23.5,2.4) node[below=3pt] {} node[above=3pt] {永};                                                          
   \draw (23.5,2.) node[below=3pt] {} node[above=3pt] {世};                                                          

 %
      \draw (10,3.8) -- (10,5); 
      \draw (12,3.8) -- (12,5); 
\draw[>=triangle 45, <->] (10,5) -- (12,5); 
  \draw (10.5,5.0) node[below=3pt] {} node[above=3pt] {教};                                                
  \draw (10.9,5.0) node[below=3pt] {} node[above=3pt] {会};                                                
    \draw (11.3,5.0) node[below=3pt] {} node[above=3pt] {时};                                                          
   \draw (11.7,5.0) node[below=3pt] {} node[above=3pt] {期};                                                          


      \draw (15,4) -- (15,5); 
\draw[>=triangle 45, <->] (12,5) -- (15,5); 
  \draw (13.,5.0) node[below=3pt] {} node[above=3pt] {圣};                                                
  \draw (13.4,5.0) node[below=3pt] {} node[above=3pt] {徒};                                                
    \draw (13.8,5.0) node[below=3pt] {} node[above=3pt] {得};                                                          
   \draw (14.2,5.0) node[below=3pt] {} node[above=3pt] {赏};                                                          
  \draw (13.,4.36) node[below=3pt] {} node[above=3pt] {地};                                                
  \draw (13.4,4.36) node[below=3pt] {} node[above=3pt] {上};                                                
    \draw (13.8,4.36) node[below=3pt] {} node[above=3pt] {灾};                                                          
   \draw (14.2,4.35) node[below=3pt] {} node[above=3pt] {难};                                                          


       \draw (-0.5,2) -- (25,2); 
      % \draw (0.,2.1) -- (0,2);       
        \draw (0.5,2.1) -- (0.5,2);    
      % \draw (1.,2.1) -- (1,2);       
        \draw (1.5,2.1) -- (1.5,2);    
      % \draw (2.,2.1) -- (2,2);       
        \draw (3.5,2.1) -- (3.5,2);   
       \draw (16,2.1) -- (16.,2);     
       \draw (17.,2.1) -- (17.,2); 
     \draw (18.,2.1) -- (18.,2); 
       \draw (19.,2.1) -- (19.,2);           
       \draw (21.5,2.1) -- (21.5,2); 
      \draw (22.5,2.1) -- (22.5,2);  
      \draw (23.5,2.1) -- (23.5,2);  
     \iffalse      
      \draw (23.8,2.1) -- (23.8,2);
     \draw (24.8,2.1) -- (24.8,2);
     \draw (25.8,2.1) -- (25.8,2);    
      \fi                                                                   
      \fill (2.5,2)  circle[radius=3pt];
   \fill (6.,2)  circle[radius=3pt]; 
   \fill (15.,2)  circle[radius=3pt];                                           
     \fill (19.2,2)  circle[radius=3pt]; 
     \fill (23.5,2)  circle[radius=3pt];  
\end{tikzpicture}
\end{sidewaystable} 
%\end{changemargin}


